\begin{table}[t]
\centering
\small
\setlength{\tabcolsep}{3.0pt}
\renewcommand{\arraystretch}{1.05}
\resizebox{\columnwidth}{!}{%
\begin{tabular}{@{}llccc@{}}
\toprule
\textbf{Robot} & \textbf{Author} & \textbf{Program length} & \textbf{Attributes/program} & \textbf{Types/Cue} \\
\midrule
IIWA & Human & $2.7 \pm 1.1$ & $10.3 \pm 4.9$ & $0.77$ \\
 & LLM & $3.0 \pm 0.7$ & $11.4 \pm 2.4$ & $0.83$ \\
\addlinespace[1pt]
Stretch & Human & $3.9 \pm 1.6$ & $15.3 \pm 6.5$ & $1.33$ \\
 & LLM & $2.8 \pm 0.6$ & $10.7 \pm 1.9$ & $1.10$ \\
\addlinespace[1pt]
Roomba & Human & $2.8 \pm 1.7$ & $11.2 \pm 6.5$ & $1.57$ \\
 & LLM & $3.2 \pm 0.9$ & $9.1 \pm 3.7$ & $1.20$ \\
\bottomrule
\end{tabular}
}
\caption{Structural complexity and compositional diversity of matched human- and LLM-authored Motion Programs ($n=30$ cues per row; mean $\pm$ SD). Program length counts top-level primitives. Attributes/program sums all explicitly written positional and named arguments; parser-inserted defaults are excluded. Types/Cue is the mean number of distinct dynamic motion families---oscillation, line, arc, and one-way joint motion---used in one program; static \texttt{pose} and \texttt{hold} primitives are excluded.}
\label{tab:nonanthro-human-llm-30}
\end{table}
